\SourceLeaf{154}
Dass ich übrigens die in der Theoria Motus Corporum Coelestium angewandte Metaphysik für die Methode der kleinsten Quadrate späterhin habe fallen lassen, ist vorzugsweise auch aus einem Grunde geschehen, den ich selbst öffentlich nicht erwähnt habe. Ich muss es nemlich in alle Wege für weniger wichtig halten, denjenigen Werth einer unbekannten Grösse auszumitteln, dessen Wahrscheinlichkeit die grösste ist, die ja doch immer nur unendlich klein bleibt, als vielmehr denjenigen, an welchen sich haltend man das am wenigsten nachtheilige Spiel hat; oder wenn $fa$ die Wahrscheinlichkeit des Werths $a$ für die Unbekannte $x$ bezeichnet, so ist weniger daran gelegen, dass $fa$ ein Maximum werde, als daran, dass $\int fx \cdot F(x-a)\,dx$, ausgedehnt durch alle möglichen Werthe des $x$, ein Minimum werde, indem für $F$ eine Function gewählt wird, die immer positiv und für grössere Argumente auf eine schickliche Art immer grösser wird. Dass man dafür das Quadrat wählt, ist rein willkürlich und diese Willkürlichkeit liegt in der Natur der Sache. Ohne die bekannten ausserordentlich grossen Vortheile, die die Wahl des Quadrats gewährt, könnte man jede andere jenen Bedingungen entsprechende wählen und thut es auch in ganz singulären Fällen. Ich weiss aber nicht, ob ich mich hinlänglich klar über das, was ich eigentlich mit dieser Distinction meine, ausgedrückt habe, und bitte im entgegengesetzten Fall um gütige Entschuldigung.

. . . . . . . . . . . . . . . .

\Needspace{6\baselineskip}
\begin{center}\textsc{Gauss} an \textsc{Schumacher}. Göttingen, 25. November 1844.\end{center}

. . . . . . . . . . . .

Beim mündlichen Vortrage der Lehre von der Methode der kleinsten Quadrate pflege ich gerade den umgekehrten Weg von demjenigen zu nehmen, den in einer gedruckten Abhandlung einzuschlagen verständig ist. Ich lehre nemlich \textls{zuerst} die Art sie anzuwenden, nach Massgabe der Umstände, mehr oder weniger von den feinern Kunstgriffen einmischend. Dann erst nehme ich, soweit dazu Zeit übrig bleibt, die verschiedenen Begründungsarten vor,
